\section{}

\paragraph{1206年}\eventref{YUAN-1206-EVENT-151}{YUAN-1206-EVENT-151}　铁木真获成吉思汗称号，蒙古诸部统一。此项以《元史》〈本纪〉及相关志；《元朝秘史》；《元典章》或《通制条格》相关条 为基本来源线索；编年与官修材料能够确认年序和制度用语，但对地方执行、人数、动机或社会后果的判断仍须同文书、碑刻、考古及对方叙事互校。

\paragraph{1206年}\eventanchor{YUAN-1206-REGNAL-YEAR}{YUAN-1206-REGNAL-YEAR}　蒙古以成吉思汗在位第1年作为诏令、赋役与外交文书的纪年框架。此项以《元史》〈本纪〉及相关志；《元朝秘史》；《元典章》或《通制条格》相关条 为基本来源线索；编年与官修材料能够确认年序和制度用语，但对地方执行、人数、动机或社会后果的判断仍须同文书、碑刻、考古及对方叙事互校。

\paragraph{1207年}\eventanchor{YUAN-1207-REGNAL-YEAR}{YUAN-1207-REGNAL-YEAR}　蒙古以成吉思汗在位第2年作为诏令、赋役与外交文书的纪年框架。此项以《元史》〈本纪〉及相关志；《元朝秘史》；《元典章》或《通制条格》相关条 为基本来源线索；编年与官修材料能够确认年序和制度用语，但对地方执行、人数、动机或社会后果的判断仍须同文书、碑刻、考古及对方叙事互校。

\paragraph{1208年}\eventanchor{YUAN-1208-REGNAL-YEAR}{YUAN-1208-REGNAL-YEAR}　蒙古以成吉思汗在位第3年作为诏令、赋役与外交文书的纪年框架。此项以《元史》〈本纪〉及相关志；《元朝秘史》；《元典章》或《通制条格》相关条 为基本来源线索；编年与官修材料能够确认年序和制度用语，但对地方执行、人数、动机或社会后果的判断仍须同文书、碑刻、考古及对方叙事互校。

\paragraph{1209年}\eventanchor{YUAN-1209-EVENT-152}{YUAN-1209-EVENT-152}　蒙古军攻西夏并迫使其接受和议与贡赋。此项以《元史》〈本纪〉及相关志；《元朝秘史》；《元典章》或《通制条格》相关条 为基本来源线索；编年与官修材料能够确认年序和制度用语，但对地方执行、人数、动机或社会后果的判断仍须同文书、碑刻、考古及对方叙事互校。

\paragraph{1209年}\eventanchor{YUAN-1209-REGNAL-YEAR}{YUAN-1209-REGNAL-YEAR}　蒙古以成吉思汗在位第4年作为诏令、赋役与外交文书的纪年框架。此项以《元史》〈本纪〉及相关志；《元朝秘史》；《元典章》或《通制条格》相关条 为基本来源线索；编年与官修材料能够确认年序和制度用语，但对地方执行、人数、动机或社会后果的判断仍须同文书、碑刻、考古及对方叙事互校。

\paragraph{1210年}\eventanchor{YUAN-1210-REGNAL-YEAR}{YUAN-1210-REGNAL-YEAR}　蒙古以成吉思汗在位第5年作为诏令、赋役与外交文书的纪年框架。此项以《元史》〈本纪〉及相关志；《元朝秘史》；《元典章》或《通制条格》相关条 为基本来源线索；编年与官修材料能够确认年序和制度用语，但对地方执行、人数、动机或社会后果的判断仍须同文书、碑刻、考古及对方叙事互校。

\paragraph{1211年}\eventanchor{YUAN-1211-EVENT-153}{YUAN-1211-EVENT-153}　蒙古大举进攻金朝北方边境。此项以《元史》〈本纪〉及相关志；《元朝秘史》；《元典章》或《通制条格》相关条 为基本来源线索；编年与官修材料能够确认年序和制度用语，但对地方执行、人数、动机或社会后果的判断仍须同文书、碑刻、考古及对方叙事互校。

\paragraph{1211年}\eventanchor{YUAN-1211-REGNAL-YEAR}{YUAN-1211-REGNAL-YEAR}　蒙古以成吉思汗在位第6年作为诏令、赋役与外交文书的纪年框架。此项以《元史》〈本纪〉及相关志；《元朝秘史》；《元典章》或《通制条格》相关条 为基本来源线索；编年与官修材料能够确认年序和制度用语，但对地方执行、人数、动机或社会后果的判断仍须同文书、碑刻、考古及对方叙事互校。

\paragraph{1212年}\eventanchor{YUAN-1212-REGNAL-YEAR}{YUAN-1212-REGNAL-YEAR}　蒙古以成吉思汗在位第7年作为诏令、赋役与外交文书的纪年框架。此项以《元史》〈本纪〉及相关志；《元朝秘史》；《元典章》或《通制条格》相关条 为基本来源线索；编年与官修材料能够确认年序和制度用语，但对地方执行、人数、动机或社会后果的判断仍须同文书、碑刻、考古及对方叙事互校。

\paragraph{1213年}\eventanchor{YUAN-1213-REGNAL-YEAR}{YUAN-1213-REGNAL-YEAR}　蒙古以成吉思汗在位第8年作为诏令、赋役与外交文书的纪年框架。此项以《元史》〈本纪〉及相关志；《元朝秘史》；《元典章》或《通制条格》相关条 为基本来源线索；编年与官修材料能够确认年序和制度用语，但对地方执行、人数、动机或社会后果的判断仍须同文书、碑刻、考古及对方叙事互校。

\paragraph{1214年}\eventanchor{YUAN-1214-REGNAL-YEAR}{YUAN-1214-REGNAL-YEAR}　蒙古以成吉思汗在位第9年作为诏令、赋役与外交文书的纪年框架。此项以《元史》〈本纪〉及相关志；《元朝秘史》；《元典章》或《通制条格》相关条 为基本来源线索；编年与官修材料能够确认年序和制度用语，但对地方执行、人数、动机或社会后果的判断仍须同文书、碑刻、考古及对方叙事互校。

\paragraph{1215年}\eventanchor{YUAN-1215-EVENT-154}{YUAN-1215-EVENT-154}　蒙古军攻占中都。此项以《元史》〈本纪〉及相关志；《元朝秘史》；《元典章》或《通制条格》相关条 为基本来源线索；编年与官修材料能够确认年序和制度用语，但对地方执行、人数、动机或社会后果的判断仍须同文书、碑刻、考古及对方叙事互校。

\paragraph{1215年}\eventanchor{YUAN-1215-REGNAL-YEAR}{YUAN-1215-REGNAL-YEAR}　蒙古以成吉思汗在位第10年作为诏令、赋役与外交文书的纪年框架。此项以《元史》〈本纪〉及相关志；《元朝秘史》；《元典章》或《通制条格》相关条 为基本来源线索；编年与官修材料能够确认年序和制度用语，但对地方执行、人数、动机或社会后果的判断仍须同文书、碑刻、考古及对方叙事互校。

\paragraph{1216年}\eventanchor{YUAN-1216-REGNAL-YEAR}{YUAN-1216-REGNAL-YEAR}　蒙古以成吉思汗在位第11年作为诏令、赋役与外交文书的纪年框架。此项以《元史》〈本纪〉及相关志；《元朝秘史》；《元典章》或《通制条格》相关条 为基本来源线索；编年与官修材料能够确认年序和制度用语，但对地方执行、人数、动机或社会后果的判断仍须同文书、碑刻、考古及对方叙事互校。

\paragraph{1217年}\eventanchor{YUAN-1217-REGNAL-YEAR}{YUAN-1217-REGNAL-YEAR}　蒙古以成吉思汗在位第12年作为诏令、赋役与外交文书的纪年框架。此项以《元史》〈本纪〉及相关志；《元朝秘史》；《元典章》或《通制条格》相关条 为基本来源线索；编年与官修材料能够确认年序和制度用语，但对地方执行、人数、动机或社会后果的判断仍须同文书、碑刻、考古及对方叙事互校。

\paragraph{1218年}\eventanchor{YUAN-1218-EVENT-155}{YUAN-1218-EVENT-155}　蒙古消灭西辽并进入中亚政治网络。此项以《元史》〈本纪〉及相关志；《元朝秘史》；《元典章》或《通制条格》相关条 为基本来源线索；编年与官修材料能够确认年序和制度用语，但对地方执行、人数、动机或社会后果的判断仍须同文书、碑刻、考古及对方叙事互校。

\paragraph{1218年}\eventanchor{YUAN-1218-REGNAL-YEAR}{YUAN-1218-REGNAL-YEAR}　蒙古以成吉思汗在位第13年作为诏令、赋役与外交文书的纪年框架。此项以《元史》〈本纪〉及相关志；《元朝秘史》；《元典章》或《通制条格》相关条 为基本来源线索；编年与官修材料能够确认年序和制度用语，但对地方执行、人数、动机或社会后果的判断仍须同文书、碑刻、考古及对方叙事互校。

\paragraph{1219年}\eventanchor{YUAN-1219-EVENT-156}{YUAN-1219-EVENT-156}　蒙古发动对花剌子模的大规模战争。此项以《元史》〈本纪〉及相关志；《元朝秘史》；《元典章》或《通制条格》相关条 为基本来源线索；编年与官修材料能够确认年序和制度用语，但对地方执行、人数、动机或社会后果的判断仍须同文书、碑刻、考古及对方叙事互校。

\paragraph{1219年}\eventanchor{YUAN-1219-REGNAL-YEAR}{YUAN-1219-REGNAL-YEAR}　蒙古以成吉思汗在位第14年作为诏令、赋役与外交文书的纪年框架。此项以《元史》〈本纪〉及相关志；《元朝秘史》；《元典章》或《通制条格》相关条 为基本来源线索；编年与官修材料能够确认年序和制度用语，但对地方执行、人数、动机或社会后果的判断仍须同文书、碑刻、考古及对方叙事互校。

\paragraph{1220年}\eventanchor{YUAN-1220-REGNAL-YEAR}{YUAN-1220-REGNAL-YEAR}　蒙古以成吉思汗在位第15年作为诏令、赋役与外交文书的纪年框架。此项以《元史》〈本纪〉及相关志；《元朝秘史》；《元典章》或《通制条格》相关条 为基本来源线索；编年与官修材料能够确认年序和制度用语，但对地方执行、人数、动机或社会后果的判断仍须同文书、碑刻、考古及对方叙事互校。

\paragraph{1221年}\eventanchor{YUAN-1221-REGNAL-YEAR}{YUAN-1221-REGNAL-YEAR}　蒙古以成吉思汗在位第16年作为诏令、赋役与外交文书的纪年框架。此项以《元史》〈本纪〉及相关志；《元朝秘史》；《元典章》或《通制条格》相关条 为基本来源线索；编年与官修材料能够确认年序和制度用语，但对地方执行、人数、动机或社会后果的判断仍须同文书、碑刻、考古及对方叙事互校。

\paragraph{1222年}\eventanchor{YUAN-1222-REGNAL-YEAR}{YUAN-1222-REGNAL-YEAR}　蒙古以成吉思汗在位第17年作为诏令、赋役与外交文书的纪年框架。此项以《元史》〈本纪〉及相关志；《元朝秘史》；《元典章》或《通制条格》相关条 为基本来源线索；编年与官修材料能够确认年序和制度用语，但对地方执行、人数、动机或社会后果的判断仍须同文书、碑刻、考古及对方叙事互校。

\paragraph{1223年}\eventanchor{YUAN-1223-REGNAL-YEAR}{YUAN-1223-REGNAL-YEAR}　蒙古以成吉思汗在位第18年作为诏令、赋役与外交文书的纪年框架。此项以《元史》〈本纪〉及相关志；《元朝秘史》；《元典章》或《通制条格》相关条 为基本来源线索；编年与官修材料能够确认年序和制度用语，但对地方执行、人数、动机或社会后果的判断仍须同文书、碑刻、考古及对方叙事互校。

\paragraph{1224年}\eventanchor{YUAN-1224-REGNAL-YEAR}{YUAN-1224-REGNAL-YEAR}　蒙古以成吉思汗在位第19年作为诏令、赋役与外交文书的纪年框架。此项以《元史》〈本纪〉及相关志；《元朝秘史》；《元典章》或《通制条格》相关条 为基本来源线索；编年与官修材料能够确认年序和制度用语，但对地方执行、人数、动机或社会后果的判断仍须同文书、碑刻、考古及对方叙事互校。

\paragraph{1225年}\eventanchor{YUAN-1225-REGNAL-YEAR}{YUAN-1225-REGNAL-YEAR}　蒙古以成吉思汗在位第20年作为诏令、赋役与外交文书的纪年框架。此项以《元史》〈本纪〉及相关志；《元朝秘史》；《元典章》或《通制条格》相关条 为基本来源线索；编年与官修材料能够确认年序和制度用语，但对地方执行、人数、动机或社会后果的判断仍须同文书、碑刻、考古及对方叙事互校。

\paragraph{1226年}\eventanchor{YUAN-1226-REGNAL-YEAR}{YUAN-1226-REGNAL-YEAR}　蒙古以成吉思汗在位第21年作为诏令、赋役与外交文书的纪年框架。此项以《元史》〈本纪〉及相关志；《元朝秘史》；《元典章》或《通制条格》相关条 为基本来源线索；编年与官修材料能够确认年序和制度用语，但对地方执行、人数、动机或社会后果的判断仍须同文书、碑刻、考古及对方叙事互校。

\paragraph{1227年}\eventanchor{YUAN-1227-EVENT-157}{YUAN-1227-EVENT-157}　西夏覆亡，同年成吉思汗去世。此项以《元史》〈本纪〉及相关志；《元朝秘史》；《元典章》或《通制条格》相关条 为基本来源线索；编年与官修材料能够确认年序和制度用语，但对地方执行、人数、动机或社会后果的判断仍须同文书、碑刻、考古及对方叙事互校。

\paragraph{1227年}\eventanchor{YUAN-1227-REGNAL-YEAR}{YUAN-1227-REGNAL-YEAR}　蒙古以成吉思汗在位第22年作为诏令、赋役与外交文书的纪年框架。此项以《元史》〈本纪〉及相关志；《元朝秘史》；《元典章》或《通制条格》相关条 为基本来源线索；编年与官修材料能够确认年序和制度用语，但对地方执行、人数、动机或社会后果的判断仍须同文书、碑刻、考古及对方叙事互校。

\paragraph{1228年}\eventanchor{YUAN-1228-REGNAL-YEAR}{YUAN-1228-REGNAL-YEAR}　蒙古以窝阔台在位第1年作为诏令、赋役与外交文书的纪年框架。此项以《元史》〈本纪〉及相关志；《元朝秘史》；《元典章》或《通制条格》相关条 为基本来源线索；编年与官修材料能够确认年序和制度用语，但对地方执行、人数、动机或社会后果的判断仍须同文书、碑刻、考古及对方叙事互校。

\paragraph{1229年}\eventanchor{YUAN-1229-EVENT-158}{YUAN-1229-EVENT-158}　窝阔台即大汗位。此项以《元史》〈本纪〉及相关志；《元朝秘史》；《元典章》或《通制条格》相关条 为基本来源线索；编年与官修材料能够确认年序和制度用语，但对地方执行、人数、动机或社会后果的判断仍须同文书、碑刻、考古及对方叙事互校。

\paragraph{1229年}\eventanchor{YUAN-1229-REGNAL-YEAR}{YUAN-1229-REGNAL-YEAR}　蒙古以窝阔台在位第2年作为诏令、赋役与外交文书的纪年框架。此项以《元史》〈本纪〉及相关志；《元朝秘史》；《元典章》或《通制条格》相关条 为基本来源线索；编年与官修材料能够确认年序和制度用语，但对地方执行、人数、动机或社会后果的判断仍须同文书、碑刻、考古及对方叙事互校。

\paragraph{1230年}\eventanchor{YUAN-1230-REGNAL-YEAR}{YUAN-1230-REGNAL-YEAR}　蒙古以窝阔台在位第3年作为诏令、赋役与外交文书的纪年框架。此项以《元史》〈本纪〉及相关志；《元朝秘史》；《元典章》或《通制条格》相关条 为基本来源线索；编年与官修材料能够确认年序和制度用语，但对地方执行、人数、动机或社会后果的判断仍须同文书、碑刻、考古及对方叙事互校。

\paragraph{1231年}\eventanchor{YUAN-1231-EVENT-159}{YUAN-1231-EVENT-159}　三峰山等战事削弱金军主力。此项以《元史》〈本纪〉及相关志；《元朝秘史》；《元典章》或《通制条格》相关条 为基本来源线索；编年与官修材料能够确认年序和制度用语，但对地方执行、人数、动机或社会后果的判断仍须同文书、碑刻、考古及对方叙事互校。

\paragraph{1231年}\eventanchor{YUAN-1231-REGNAL-YEAR}{YUAN-1231-REGNAL-YEAR}　蒙古以窝阔台在位第4年作为诏令、赋役与外交文书的纪年框架。此项以《元史》〈本纪〉及相关志；《元朝秘史》；《元典章》或《通制条格》相关条 为基本来源线索；编年与官修材料能够确认年序和制度用语，但对地方执行、人数、动机或社会后果的判断仍须同文书、碑刻、考古及对方叙事互校。

